\begin{scanpage}{0069}{52}
they are interpolated. Since the transition from Equations (8.12) to Equations (8.17) is not completely unique, in the smaller interval $0<\theta\leq\frac1{ka}$ Equations (8.17) may give a certain error. This error is not very significant, since in this interval the field from the uniform part of the current is large.
\sectiontitle{9}{The Total Field Being Scattered by a Disk with Normal Irradiation}
Turning to Equations (8.07), (8.08) and (8.17), let us represent the field from the nonuniform part of the current in the following form:
\[\left.\begin{aligned}E_\phi=-H_\theta&=iaH_{0x}\frac{\cos\phi}{\sin\theta}J_1(ka\sin\theta)\frac{e^{ikR}}R-\\&\quad-\frac{iaH_{0x}}2\{[f(2)-f(1)]J_1(ka\sin\theta)+\\&\qquad+i[f(2)+f(1)]J_2(ka\sin\theta)\}\cos\phi\frac{e^{ikR}}R,\\E_\theta=H_\phi&=iaH_{0x}\frac{\sin\phi}{\sin\theta}\cos\theta J_1(ka\sin\theta)\frac{e^{ikR}}R-\\&\quad-\frac{iaH_{0x}}2\{[g(2)-g(1)]J_1(ka\sin\theta)+\\&\qquad+i[g(2)+g(1)]J_2(ka\sin\theta)\}\sin\phi\frac{e^{ikR}}R.\end{aligned}\right\}\tag{9.01}\]
Here the first terms, as is readily apparent, represent the field from the uniform part of the current taken with the opposite sign. As a result, the total field scattered by the disk (that is, the sum of the fields radiated by the uniform and nonuniform parts of the current) will be expressed only in terms of the functions $f$ and $g$ which determine in the rigorous solution the cylindrical wave from the half-plane's edge
\[\left.\begin{aligned}E_\phi=-H_\theta&=-\frac{iaH_{0x}}2\cos\phi\{[f(2)-f(1)]J_1(ka\sin\theta)+\\&\qquad+i[f(2)+f(1)]J_2(ka\sin\theta)\}\frac{e^{ikR}}R,\\E_\theta=H_\phi&=-\frac{iaH_{0x}}2\sin\phi\{[g(2)-g(1)]J_1(ka\sin\theta)+\\&\qquad+i[g(2)+g(1)]J_2(ka\sin\theta)\}\frac{e^{ikR}}R.\end{aligned}\right\}\tag{9.02}\]
\end{scanpage}
